\documentclass[11pt]{article}

% ACL 2027 submission: keep [review] for double-blind submission.
\usepackage[review]{acl}

% Standard package includes
\usepackage{times}
\usepackage{latexsym}
\usepackage[T1]{fontenc}
\usepackage[utf8]{inputenc}
\usepackage{microtype}
\usepackage{inconsolata}
\usepackage{graphicx}
\usepackage{booktabs}
\usepackage{amsmath}

\title{GUI State Recovery: Auditing and Repairing Environment Contamination Left by Mobile GUI Agents}

\author{Anonymous ACL submission}

\begin{document}
\maketitle

\begin{abstract}
% TODO: draft abstract. Key points to cover:
% (1) Mobile GUI agents leave persistent state changes (contamination)
%     after task execution, e.g., search histories and settings toggles,
%     which corrupt subsequent evaluations and real-user environments.
% (2) We formalize GUI state recovery and propose a closed-loop framework:
%     Audit Agent (infers a contamination list from execution trajectories),
%     Planner Agent (generates recovery tasks), Execute Agent (repairs the
%     state on real devices), with recovery trajectories re-audited until
%     the state is clean or escalated to a human gate.
% (3) Results on MobileBench-OL (310 tasks, 310 human golden labels) and
%     the MBOL+ extension (859 tasks); end-to-end recovery success rate
%     and inference precision/recall/F1 (all figures traceable to human
%     annotations only).
\end{abstract}

\section{Introduction}
% TODO: motivation, gap, contributions (3 bullets).

\section{Related Work}
% TODO: mobile GUI agents and benchmarks; environment state and
% reproducibility; memory/state editing and environment repair.

\section{Task Formulation}
% TODO: define contamination, contamination list schema
% (contamination-list/v0.1), recovery task schema (recovery-task/v0.1),
% success criterion, and the reset-free setting.

\section{Method}
% TODO: overall closed-loop figure (contamination monitoring ->
% recovery task generation -> execution -> re-audit).
\subsection{Audit Agent: Contamination Inference}
% TODO: deterministic ledger carry-over, search-term normalization,
% bidirectional grounding, three deterministic guards with human gating.
\subsection{Planner Agent: Recovery Task Generation}
% TODO: four-layer constraint architecture: skeleton pre-fill,
% per-mutation generation, prompt constraints, structured output validation.
\subsection{Execute Agent and Recovery Harness}
% TODO: benchmark-identical action space (7 actions), action contract,
% three-way post-action verification, Policy Engine, dynamic control loop.
\subsection{Closed-Loop Verification}
% TODO: VLM verification anchored on contamination entities; two-stage
% verification (task completion vs. contamination removal).

\section{Experimental Setup}
% TODO: datasets (MobileBench-OL base 310; MBOL+ 859), models
% (qwen3-vl-plus across all stages), human golden-label annotation
% protocol and quality checks, evaluation metrics (P, R, strict F1,
% type-only F1; recovery success rate, one-shot rate, iterations).

\section{Results}
% TODO: tables only with human-traceable numbers.
\subsection{Contamination Inference Quality}
% TODO: modality ablation (text / +XML / +image / +XML+image) for the two
% visually capable models.
\subsection{End-to-End Recovery}
% TODO: full-run statistics and failure attribution (system vs. model).

\section{Analysis}
% TODO: case studies, failure modes (model over-conservatism on
% irreversible operations, repeated plans), iteration cost.

\section{Limitations}

\section*{Ethical Statement}

\section{Conclusion}

\bibliography{custom}

\end{document}
